% =============================================================================
% ARC Solvers: Mechanistic Through Lines
% A Synthesis of URM, TRM, HRM, and VARC
% =============================================================================
\documentclass[11pt, letterpaper]{article}

% --- Packages ---
\usepackage[margin=1in]{geometry}
\usepackage{booktabs}
\usepackage{amsmath}
\usepackage{hyperref}
\usepackage{xcolor}
\usepackage{float}
\usepackage{caption}
\usepackage[utf8]{inputenc}
\usepackage[T1]{fontenc}
\usepackage{authblk}
\usepackage{url}
\usepackage{parskip}
\usepackage[numbers,sort&compress]{natbib}

% --- Colors ---
\definecolor{darkslate}{HTML}{2C3E50}
\definecolor{thinkingteal}{HTML}{1ABC9C}
\definecolor{alertred}{HTML}{E31A1C}

% --- Hyperlinks ---
\hypersetup{
    colorlinks=true,
    linkcolor=darkslate,
    citecolor=darkslate,
    urlcolor=darkslate,
}

% --- Metadata ---
\title{\textbf{Mechanistic Through Lines in ARC Solvers:\\
Iterative Computation, Inductive Bias, and Test-Time Search}}

\author{@notcomplex\_}
\date{March 2026}

% =============================================================================
\begin{document}
\maketitle

% =============================================================================
\begin{abstract}
Four recent ARC-solving architectures---Universal Reasoning Model (URM)
\cite{gao2025urm}, Vision ARC (VARC) \cite{hu2025varc}, Tiny Recursive Model
(TRM) \cite{jolicoeur2025trm}, and McGovern's test-time adaptation study
\cite{mcgovern2025ttatrm}---differ in apparent motivation yet converge on the
same underlying computational strategy: \emph{iterate a hypothesis, score it
against demonstrations, select among candidates, and squeeze the search space
with strong inductive bias}. This paper synthesizes those four works into a
single mechanistic account, identifies the load-bearing contributions that
explain performance, and surfaces the deeper open question that none of the
papers fully answers: \emph{what is the loop actually computing?}
\end{abstract}

% =============================================================================
\section{What ARC Actually Demands}
\label{sec:what_arc_demands}

The Abstraction and Reasoning Corpus (ARC) was designed as a testbed for
\emph{fluid intelligence}: the ability to identify and apply a rule from only
a handful of input-output examples, in a domain never encountered
before~\cite{chollet2019arc}. The benchmark's structure is deliberately
hostile to standard machine learning: there are roughly 400 training tasks,
each presented as two to four demonstration pairs and one held-out test
input, where the test target is unknown. No task appears in two splits.
ARC-AGI-2 tightens that screws further by calibrating tasks to be harder
for symbolic brute-force while remaining tractable to humans~\cite{arcprizefoundation2025}.

Two structural facts about ARC shape everything that follows:

\textbf{The sample budget is tiny.} Two to four examples is not enough to fit
even a small neural network from scratch. Any successful system must either
(a) carry strong priors capable of extrapolating from near-zero data, or (b)
use the demonstrations as an on-line search oracle rather than a training set.

\textbf{Evaluation is pixel-perfect.} A single incorrect cell anywhere in the
output grid fails the task. This makes probabilistic averaging---averaging
multiple distinct hypotheses pixel-by-pixel---almost always harmful, and
therefore forces systems toward discrete hypothesis selection rather than
soft ensemble output.

Together these two constraints define the design space that URM, TRM, and
VARC navigate in distinct but compatible ways.

% =============================================================================
\section{The Unifying Loop}
\label{sec:loop}

Each of the four architectures reviewed here can be described as an instance
of the same computational skeleton:

\begin{enumerate}
  \item \textbf{Encode.} Represent the task (demonstrations + query) as a
        structured embedding or token sequence.
  \item \textbf{Refine.} Apply a fixed update operator repeatedly: either a
        shared-weight recurrent block (URM, TRM, HRM) or a gradient update
        on a task token (VARC).
  \item \textbf{Decode.} Read off a candidate output.
  \item \textbf{Select.} Choose among many candidate outputs by consistency
        with the demonstrations (voting, majority, or program validation).
\end{enumerate}

The loop is not a metaphor. In URM/TRM/HRM it is literally a recurrent unrolled
computation; in VARC it is gradient descent applied per-task at inference time.
The convergence of these superficially different approaches on the same four-step
skeleton is the central observation this paper attempts to explain.

% =============================================================================
\section{Looped Transformers: URM and TRM}
\label{sec:looped}

\subsection{The core claim: recurrence beats depth}

The Universal Transformer (UT)~\cite{dehghani2019ut} replaced a stack of
distinct layers with repeated application of a single shared-weight block.
This is architecturally embarrassingly simple, but URM's ablations provide
the clearest evidence in the literature that the simplicity is doing real
work~\cite{gao2025urm}.

URM's headline comparison: at an equal 32$\times$ compute multiplier,
reallocating computation from non-shared to recurrent passes improves
pass@1 on ARC-AGI-1 from 23.75\% to 40.0\%. More importantly, simply
scaling a vanilla (non-recurrent) Transformer's depth and width under the
same budget yields only diminishing returns or degradation. This is strong
evidence for the hypothesis that \emph{iterative refinement is a better use
of compute than additional distinct parameters for this class of task}.

Why might that be? The intuition is that ARC tasks often require a sequence
of logically dependent sub-steps---detect a repeating motif, identify which
axis of symmetry applies, infer the color mapping, predict the output grid.
A fixed-depth feedforward network must commit to a single parse of the task
in one pass. A looped model can ``circle back'': an early pass produces
a rough hypothesis; later passes refine, correct, and fill in detail, with
the weight tying ensuring that the refinement operator remains consistent
across steps. The result is effective depth without parameter explosion.

\subsection{What the ablations actually isolate}

URM's ablation study goes further than any prior work in this family by
decomposing the gain into attributable components~\cite{gao2025urm}:

\textbf{Nonlinearity is load-bearing, not decorative.} Removing or weakening
nonlinear components produces monotonic degradation. Replacing SwiGLU with
SiLU or ReLU drops accuracy substantially; removing the attention softmax
collapses performance to near-random ($\sim$2\% pass@1 in their table).
This implicates the MLP block---specifically its gating nonlinearity---as the
primary source of representational capacity for ARC-style abstract logic,
not the attention head itself. Attention may be doing ``routing'' while the
MLP is doing ``thinking.''

\textbf{ConvSwiGLU: convolution in the gate.} URM inserts a short depthwise
convolution into the MLP gating pathway. Their diagnostic is that placing
the convolution \emph{after} expansion is most effective, consistent with
the interpretation that it is augmenting the expressive capacity of the
nonlinear mixing stage rather than preprocessing features before gating.
The practical gain is meaningful: this single addition substantially closes
the gap to VARC on ARC-AGI-1.

\textbf{Truncated Backpropagation Through Loops (TBPTL).} Long unrolled
recurrences are hard to train: gradients either vanish or explode as they
flow backward through many identical blocks. URM's solution---compute
gradients only through the later loops, analogous to truncated BPTT in
sequence modeling---yields a reported ``sweet spot'' at skipping the first
two inner loops. This is mechanistically interesting: the claim is that
early loops are the exploratory phase (rough hypothesis generation) and
later loops are the refinement phase (precision correction), and that useful
gradient signal lives mostly in the latter.

\subsection{TRM: the scratchpad interpretation}

TRM arrives at a cleaner mechanistic story for what the latent state
\emph{is}~\cite{jolicoeur2025trm}. Rather than two networks operating at
different biological timescales (as HRM posited), TRM proposes a single tiny
2-layer network and two latent variables with explicit roles: one latent is
the current candidate answer, the other is a reasoning scratchpad ``acting
similarly as a chain-of-thought, except in latent space.''

This reframing matters because it gives us an interpretable decomposition:
the model is simultaneously maintaining a guess and a process for improving
that guess. The scratchpad is not decoded; it is not supervised directly;
it exists purely to make the next refinement step more accurate. This is
functionally equivalent to hidden chain-of-thought, and the evidence that
it works ($\sim$45\% pass@1 on ARC-AGI-1 with only 7M parameters) suggests
that latent scratchpads may be a computationally efficient substitute for
explicit verbal reasoning traces.

TRM also makes test-time augmentation a central method rather than a
post-hoc trick: at inference the model runs across $\sim$1000 augmented
views (color permutations, dihedral transforms, translations) and reports
the plurality answer. This is not a minor implementation detail---it is
functionally a test-time search over a structured space of equivalent
task restatements, and it contributes substantially to the reported accuracy.

% =============================================================================
\section{Vision as Inductive Prior: VARC}
\label{sec:varc}

VARC's central argument is that the dominant difficulty in ARC is not
computational but representational: the wrong modality creates an
artificially hard inductive-bias mismatch~\cite{hu2025varc}.

\subsection{ARC tasks are visual at their core}

Many ARC transformations---reflection, symmetry, gravity, object
translation---are primitives of 2D visual and physical experience. They
are ``difficult'' for language models partly because language models
tokenize grids into flat sequences, discarding the 2D structure that
makes the rule legible to a human. VARC's insight is that framing each
task as image-to-image translation with a \emph{canvas} representation
(a larger spatial context in which the grids are embedded) allows standard
vision architectures---ViTs, U-Nets---to apply 2D inductive priors
directly.

The ablation evidence for this claim is unusually sharp: replacing 2D
rotary positional embeddings with 1D RoPE drops accuracy from 54.5\% to
51.0\% (a 3.5 point degradation). That is a large effect for a single
hyperparameter change, and it confirms that spatial position structure
is contributing real signal, not just acting as a cosmetic prior.

\subsection{Test-time training as task adaptation}

VARC has no recurrence within the inference pass. Once the model is trained,
inference for a new task runs as follows: initialize a task-specific token
randomly, fine-tune it (and potentially the entire model) on the few
demonstration pairs using gradient descent, then do a feedforward pass to
predict the output. This is test-time training (TTT), and it substitutes
for the recurrence loop: instead of updating a hidden state by repeatedly
applying the same block, VARC updates the model's parameters (or a task
projection) by repeatedly applying gradient descent to the demonstration
loss.

The quantitative payoff is large:
\begin{itemize}
  \item Single-view, no TTT: approximately 35.9\% pass@1
  \item Multi-view with TTT: 49.8\% pass@1
  \item Multi-view TTT, pass@2: 54.5\%
  \item Ensemble (ViT + U-Net, both with TTT): 60.4\%
\end{itemize}

The gap between the first and last row---24.5 percentage points---is a
direct measure of how much ``search'' (over views and model states) is
doing, compared to the base model capacity alone.

\subsection{Multi-view inference as hypothesis marginalization}

VARC's multi-view strategy (up to 510 augmented views, majority vote) is
structurally identical to TRM's augmentation voting. Both are marginalizing
over a transformation group: the group of geometric and color symmetries
that should leave the correct answer invariant. This is not just an
engineering trick---it is a principled Bayesian operation applied to a
finite, discrete group, and the gains it provides suggest that the base
models are sensitive to nuisance symmetries that a human solver would
automatically factor out.

% =============================================================================
\section{Test-Time Adaptation at Compute Limits: McGovern}
\label{sec:mcgovern}

McGovern's study~\cite{mcgovern2025ttatrm} is the most practically grounded
of the four: it asks whether TRM-style models can be made to work inside the
ARC Prize competition's severe compute budget (4$\times$L4 GPUs, 12 hours).

\subsection{The compute gap and what it reveals}

TRM pretraining requires approximately 750k optimizer steps at a batch size
of 768. Competition runtime permits roughly 15k steps at batch size 384.
That is a 1/32 compute ratio---a gap that cannot be closed by algorithmic
tricks alone.

McGovern's answer is to separate pretraining (unconstrained, done offline)
from competition-time fine-tuning (constrained). The finding is that a
pretrained model fine-tuned for 12,500 steps reaches 6.67\% on semi-private
tasks, compared to 10\% for the fully pretrained model. That is a 33\%
relative gap---not negligible, but also not catastrophic. The pretrained
weights provide a useful initialization.

\subsection{The task-embedding hypothesis fails}

The most mechanistically informative finding in McGovern's paper is negative:
fine-tuning only the task ID embeddings (while holding the trunk frozen)
achieves near-zero accuracy on held-out tasks. This falsifies a clean
interpretation in which the task embedding encodes the transformation rule
and the frozen trunk executes it. Substantial weight updates throughout the
network are required for any given unseen task.

This is an important constraint on interpretive stories. The ``soft program
in the embedding'' hypothesis is appealing, but the evidence suggests that
the program is distributed across the full network, not concentrated in a
dedicated code vector. Test-time adaptation must reach the trunk.

% =============================================================================
\section{What the Loop is Actually Doing: An Open Synthesis}
\label{sec:synthesis}

Taking all four papers together, the following picture emerges---though it
must be read as a mechanistic hypothesis, not an established fact.

\subsection{Two sources of reasoning, partly separable}

ARC performance appears to have two partially separable sources:

\textbf{Inductive bias / representation alignment.} The right prior shrinks
the hypothesis space dramatically before any computation is done. VARC's
2D positional encodings, canvas representation, and visual pretraining do
this work. URM's recurrent inductive bias (depth via iteration rather than
parameter count) does a different version of it. Getting the modality and
structure right is not a fine-tuning detail---it is half the problem.

\textbf{Iterative adaptation / test-time search.} Given the right prior,
each method then does a form of search: recurrence (URM/TRM) updates
a hypothesis state; TTT (VARC, McGovern) updates model weights; augmentation
voting (all four) marginalizes over nuisance symmetries. The two sources
compound: high-quality priors enable shorter, more reliable search; rich
search compensates for imperfect priors.

\subsection{The loop as latent algorithm emulation}

The theoretical framing that best fits the empirical evidence is that looped
transformers are emulating iterative algorithms in latent space. A looped
$k$-layer transformer applied $L$ times can approximate a $kL$-layer
feedforward model, but more than that, the parameter sharing across loops
induces a structure closer to a \emph{fixed-point iteration} than to a
simple deep pipeline. The loop is searching for a latent representation from
which the correct output can be read off, using the same update operator at
each step because the same type of ``correction'' is always needed---check
for consistency, refine the current hypothesis, propagate constraints.

TRM's scratchpad latent is the clearest expression of this idea: one half of
the hidden state is ``the guess,'' the other is ``the evidence for correcting
the guess.'' As loops accumulate, the guess converges and the scratchpad
empties. Whether this is analogous to human reasoning steps (object
detection $\to$ grouping $\to$ rule identification $\to$ application) or
something altogether different is unknown---none of these papers contains
mechanistic interpretability at that resolution.

\subsection{Scale enters primarily as test-time compute}

Across all four works, the gains that look like ``scale'' are actually
test-time compute: more augmentation views, more refinement loops, more
fine-tuning steps, more candidates to vote over. The parameter count of the
\emph{base model} is surprisingly small (7M for TRM, 18M for VARC). This
suggests that ARC, at least in its current form, rewards \emph{compute
efficiency under a tight parameter budget} more than raw parameter scale.
This could be a genuine signal about the nature of fluid reasoning (iterative
computation is more efficient than parametric memorization for rule
induction) or it could be an artifact of the dataset's small size making
large models overfit.

ARC-AGI-2's harder tasks and the observed gap between ARC-AGI-1 and
ARC-AGI-2 scores---URM achieves 53.8\% on V1 but only 16.0\% on
V2~\cite{gao2025urm}; TRM achieves 45\% on V1 but only 8\% on
V2~\cite{jolicoeur2025trm}---suggests that whatever these systems have
learned is real but narrow. The representations generalize within the
difficulty and concept distribution of ARC-AGI-1 but partially fail
when that distribution shifts. Whether the core mechanisms (recurrence,
visual priors, TTT) are sufficient or whether fundamentally new ideas
are required for ARC-AGI-2 remains the central open question.

% =============================================================================
\section{What Is Missing}
\label{sec:missing}

The papers reviewed here share a common pattern of evidence: ``remove the
component, observe collapse.'' URM shows this for nonlinearity and
recurrence~\cite{gao2025urm}; VARC shows it for 2D positional
encoding~\cite{hu2025varc}. This is valuable but leaves deeper questions
untouched:

\textbf{What intermediate representations arise inside the loops?} Are there
stable attractors in the latent space that correspond to interpretable
concepts (``this is a reflection task,'' ``this is a color-mapping task'')?
None of the papers contains probing experiments or representational geometry
analysis that would answer this.

\textbf{Can augmentation gains be replaced by better priors?} The large
contribution of multi-view voting (TRM: $\sim$1000 views; VARC: up to 510
views) raises the question of whether these gains reflect a genuine
marginalization over nuisance symmetries or a compensation for a
representation that is insufficiently invariant. If the latter, better
architecture might render test-time augmentation unnecessary.

\textbf{How does the test-time adaptation interact with pretraining?} McGovern
shows that full fine-tuning dominates embedding-only adaptation, but does not
decompose which layers change most or what changes about them. Understanding
this would directly inform whether larger pretrained trunks would help or
merely shift the locus of adaptation.

\textbf{Is the loop convergent?} TRM and URM both use fixed (hyperparameter-
chosen) loop counts. Neither paper presents evidence of a convergence
criterion---a condition under which the system knows it has found a good
hypothesis and should stop iterating. Adaptive halting (ACT) is available in
the UT framework but does not appear in the ARC-specific papers. Whether
adaptive depth would help or be necessary for harder tasks is an open
empirical question.

% =============================================================================
\section{Conclusion}
\label{sec:conclusion}

Four mechanistically distinct ARC solvers---URM, VARC, TRM, and the
McGovern TRM adaptation study---converge on a shared computational
strategy: \emph{iterate a hypothesis under strong inductive bias, then
select among candidates at test time}. The iteration is implemented as
recurrent weight-shared depth (URM/TRM), gradient-based task adaptation
(VARC/McGovern), and augmentation-based voting (all four). The inductive
bias varies by modality (linguistic tokens vs.\ 2D vision priors) but
serves the same function in each case: drastically shrinking the set of
plausible hypotheses before search begins.

The strongest mechanistic claim supported by actual ablation evidence is
dual: (1) iterative computation via parameter-shared recurrence is a more
efficient allocation of FLOP budget than additional distinct parameters for
multi-step abstract rule induction; and (2) the nonlinear (gating) component
of the update block is load-bearing---removing it collapses performance,
implicating ``reasoning'' in the MLP rather than the attention head.

What these systems are not yet able to do is explain their own reasoning.
The latent loops, scratchpads, and task tokens remain opaque. Closing the
gap between ``the loop converges to the right answer'' and ``the loop traces
steps recognizable as logical inference'' is, arguably, the most important
next problem in ARC mechanistic research.

% =============================================================================
\bibliographystyle{unsrtnat}
\bibliography{references}

\end{document}
